\documentclass{ppgi-resumo}
\usepackage{geometry}
\usepackage{enumitem}
\usepackage{lipsum}
\usepackage{expl3}
\usepackage[fontsize=9pt]{fontsize}
\geometry{
 a4paper,
 left=1cm,
 right=1cm,bottom=1cm,top=1.5cm,
 }

\author{
  \fcolorbox{colorUTFPR}{colorUTFPR}{\textbf{P}}rimeiro Autor\\
  \textit{autor1@utfpr.edu.br}
  \and
  \fcolorbox{colorUTFPR}{colorUTFPR}{\textbf{S}}egundo Autor\\
  \textit{autor2@utfpr.edu.br}
  \and
  \fcolorbox{colorUTFPR}{colorUTFPR}{\textbf{T}}erceiro Autor\\
  \textit{autor3@utfpr.edu.br}
}

\title{Título}


\begin{document}
   
\maketitle

\section*{Inovação}
\textcolor{colorUTFPR}{Descreva o tipo de novidade resultante de seu trabalho. Exemplos: de produto de software; de processo; em serviço; etc. Consultar \url{https://www.gov.br/capes/pt-br/centrais-de-conteudo/2020-01-03-relatorio-gt-inovacao-e-transferencia-de-conhecimento-pdf} para outras classificações de inovação}

\section*{Produção Técnica}
\textcolor{colorUTFPR}{Descreva o tipo de produção técnica resultante de seu trabalho Exemplos: software/aplicativo/programa de computador; tecnologia social; patente; material didático; etc. Consultar \url{https://www.gov.br/capes/pt-br/centrais-de-conteudo/10062019-producao-tecnica-pdf} para definição de outros tipos de produção técnica – Tabela 2, página 17}

\section*{Área de Conhecimento}
\textcolor{colorUTFPR}{São exemplos de área de conhecimento: Engenharia de Software; Banco de Dados; Processamento Gráfico; etc. Consultar 
\url{http://lattes.cnpq.br/documents/11871/24930/TabeladeAreasdoConhecimento.pdf/d192ff6b-3e0a-4074-a74d-c280521bd5f7} para definição da área de conhecimento}

\section*{Tipo de Pesquisa}
\textcolor{colorUTFPR}{Refere-se ao tipo de estudo. Utilizar classificação proposta por Wieringa et al. Requirements engineering paper classification and evaluation criteria: A proposal and a discussion.Requirements Engineering, 11(1):102-107, 2005. Exemplos de tipo de pesquisa:  pesquisa de validação; proposta de solução; relato de experiência; artigo de opinião; etc}

\section*{Objetivo}
\vspace{-0.7cm}
\noindent
    \minipage[t]{\dimexpr0.62\linewidth-0\fboxsep-0\fboxrule\relax}
        \vspace{0.4cm}
        \textcolor{colorUTFPR}{Liste o objetivo geral da pesquisa de forma direta e clara.}
        \vspace{0.3cm}
        
        \noindent\textbf{\Large{APLICAÇÃO DOS ACHADOS NA PRÁTICA}}
        \begin{itemize}[leftmargin=*]
	\setlength\itemsep{-0.3em}
	\item \textbf{[1] [}\textcolor{colorUTFPR}{Descreva de forma clara e direta as descobertas do estudo e como elas são transferíveis para configurações industriais ...}\textbf{]}
	\item \textbf{[2] [}\textcolor{colorUTFPR}{Descreva de forma clara e direta diretrizes para a prática com base nos seus resultados e se os benefícios de usar as diretrizes compensarão os custos...}\textbf{]}
	\item \textbf{[3] [}\textcolor{colorUTFPR}{etc ...}\textbf{]}
\end{itemize}
    \endminipage\hfill  
    \fcolorbox{colorUTFPR}{colorUTFPR}{ %
    \minipage[t]{\dimexpr0.35\linewidth-0\fboxsep-0\fboxrule}
        \textbf{\large PARCERIA(S)}
        \begin{itemize}[leftmargin=*]
        \setlength\itemsep{-0.2em}
        \item \textit{Cite empresas participantes}   
        \item \textit{Cite órgãos de fomento     }
        \end{itemize}
        
        \textbf{\large TÉCNICA(S) USADA(S)}
        \begin{itemize}[leftmargin=*]
        \setlength\itemsep{-0.2em}
        \item \textit{Liste o nome da técnica adotada}  
        \item \textit{Liste o nome do algoritmo usado}     
        \end{itemize}

 		\textbf{\large PALAVRAS-CHAVE}
 		\begin{itemize}[leftmargin=*]
 		\setlength\itemsep{-0.2em}
        \item \textit{Primeira palavra-chave}
        \item \textit{Segunda palavra-chave}
        \item \textit{Terceira palavra-chave}   
        \end{itemize}
    \endminipage
   }


\section*{Resumo gráfico}

\textcolor{colorUTFPR}{Crie uma imagem para ajudar o leitor a obter rapidamente uma visão geral sobre seu trabalho e entender o propósito e os resultados de sua pesquisa. Consultar \url{https://www.elsevier.com/authors/tools-and-resources/visual-abstract} para exemplos e diretrizes para elaboração de bons resumos gráficos.}

\section*{Como citar}
\begin{minipage}{\linewidth}
\color{colorUTFPR}
\begin{verbatim}
 @MastersThesis{AuthorLastNameYear,                 @PhdThesis{AuthorLastName,                            
   author = {AuthorName},                             author = {AuthorName},    
   title  = {Dissertation Title},                     title  = {Thesis Title},
   school = {University Name},                        school = {University Name},
   year   = {Year},                                   year   = {Year},
 }                                                  }
\end{verbatim}

\end{minipage}
\section*{Trabalhos Relacionados -- Leia também }

\textcolor{colorUTFPR}{Cite outras referências relacionadas à sua pesquisa, por exemplo, um artigo publicado, uma produção técnica, etc.}

\end{document}